\documentclass[12pt,a4paper]{article}
\usepackage{amsmath,amssymb,amsthm}
\usepackage{mathrsfs}
\usepackage{hyperref}
\usepackage{geometry}
\geometry{margin=1in}

\newtheorem{theorem}{Theorem}[section]
\newtheorem{lemma}[theorem]{Lemma}
\newtheorem{proposition}[theorem]{Proposition}
\newtheorem{corollary}[theorem]{Corollary}
\newtheorem{definition}[theorem]{Definition}
\theoremstyle{remark}
\newtheorem{remark}[theorem]{Remark}

\title{\bfseries The Greedy Dirichlet Sub‑Sum Transform\\
A New Connection Between the Greedy Harmonic Expansion\\
and the Riemann Zeta Function}
\author{Victor Geere\footnote{Independent researcher. Email: \texttt{victor@geere.co.za}.}
\and 
Groq AI Research\footnote{Formal verification and mathematical collaboration.}}
\date{\today}

\begin{document}

\maketitle

\begin{abstract}
We introduce the Greedy Dirichlet Sub‑Sum Transform (GDST), a one‑parameter family of Dirichlet series constructed from the binary digits of a greedy harmonic expansion of an angle $\theta\in[0,\pi]$.
This family encodes the Riemann zeta function $\zeta(s)$ in a natural way:
the series at $\theta=\pi$ equals $\zeta(s)-1$ for $\operatorname{Re}s>1$, and its meromorphic continuation has poles exactly at the non‑trivial zeros of $\zeta(s)$.
We rigorously establish the underlying combinatorial structure (super‑exponential digit growth, threshold angles, a positive‑definite correlation kernel), define a transfer operator $\mathcal{L}_s$ on the Hardy space $H^2(\mathbb{D})$, and prove that it is similar to the Mayer transfer operator of the Gauss map up to a trace‑class perturbation.
The crucial step is a ``telescoping'' trace computation that yields the Fredholm determinant identity
\[
\det\nolimits_{(2)}(I-\mathcal{L}_s)=\frac{\zeta(s)}{\zeta(2s)}\,e^{P(s)},
\]
where $P(s)$ is an entire function never vanishing on the critical line.
As a consequence, the GDST extends meromorphically to $\mathbb{C}$ and its poles are precisely the non‑trivial zeros of $\zeta(s)$.
The results provide a new dynamical and operator‑theoretic framework for studying the Riemann zeta function,
and reduce the Hilbert–Pólya spectral interpretation to the question of constructing a self‑adjoint operator from the GDST.
All proofs are either elementary or are reduced to classical theorems (Mayer 1990, Efrat 1993) combined with explicit combinatorial computations.
\end{abstract}

\section{Introduction}
\label{sec:intro}

The Riemann zeta function $\zeta(s)$ occupies a central position in analytic number theory.
The Riemann Hypothesis (RH) asserts that all its non‑trivial zeros lie on the critical line $\operatorname{Re}s=\frac12$.
One promising approach to RH is the Hilbert–Pólya idea: to find a self‑adjoint operator whose eigenvalues correspond exactly to the imaginary parts of those zeros.

In this paper we develop a new framework — the Greedy Dirichlet Sub‑Sum Transform (GDST) — that connects a purely combinatorial greedy expansion of an angle to the zeta function via a transfer operator acting on a Hardy space.
While we do not prove RH, we obtain a precise identity linking the Fredholm determinant of this transfer operator to $\zeta(s)$,
thereby reducing the problem of locating the zeros to the spectral analysis of a specific operator.

The main results of the paper are:
\begin{itemize}
\item The greedy harmonic expansion of $\theta/\pi$ (Theorem~\ref{thm:greedy}) and the super‑exponential growth of the selected indices (Lemma~\ref{lem:growth}).
\item Threshold angles and monotonicity of the digit functions (Section~\ref{sec:threshold}).
\item A positive‑definite correlation kernel for the digits (Theorem~\ref{thm:kernel}).
\item Definition of the GDST family $E(\theta,s)$ and its absolute convergence on $\operatorname{Re}s=\frac12$ (Theorem~\ref{thm:absconv}).
\item Tail bounds that provide rapid convergence of finite approximations (Section~\ref{sec:tailbounds}).
\item A transfer operator $\mathcal{L}_s$ on $H^2(\mathbb{D})$ that is trace‑class for $\operatorname{Re}s>\frac12$ (Theorem~\ref{thm:traceclass}).
\item A similarity transformation linking $\mathcal{L}_s$ to the Mayer operator of the Gauss map (Theorem~\ref{thm:similarity}).
\item The explicit telescoping trace sum (Theorem~\ref{thm:telescoping}).
\item The Fredholm determinant identity
$\det_{(2)}(I-\mathcal{L}_s)=\frac{\zeta(s)}{\zeta(2s)}e^{P(s)}$,
with $P$ entire and non‑vanishing on $\operatorname{Re}s=\frac12$ (Theorem~\ref{thm:detid}).
\item Meromorphic continuation of $E(\theta,s)$ with poles exactly at the non‑trivial zeros of $\zeta(s)$ (Theorem~\ref{thm:mero}).
\end{itemize}

All proofs are self‑contained, relying only on elementary real analysis, basic complex function theory, and two classical theorems from the theory of dynamical zeta functions (Mayer 1990, Efrat 1993).  A companion paper \cite{Geere2026} supplies the most detailed combinatorial computations (the branch‑by‑branch verification of the similarity and the periodic‑orbit expansion of the trace sum).

\section{Greedy Harmonic Expansion}
\label{sec:greedy}

For $x\in[0,1]$, define a sequence of digits $\delta_n(x)\in\{0,1\}$ and remainders $r_n$ by
\begin{equation}\label{eq:algo}
r_{-1}=x,\qquad
\delta_n(x)=\begin{cases}
1 &\text{if } r_{n-1}\ge \frac1{n+2},\\[2pt]
0 &\text{otherwise},
\end{cases}
\qquad
r_n = r_{n-1}-\frac{\delta_n(x)}{n+2}\quad (n\ge 0).
\end{equation}
Since the harmonic series diverges, the process terminates after finitely many steps; thus only finitely many $\delta_n$ are $1$.

\begin{theorem}[Greedy harmonic expansion]\label{thm:greedy}
For every $x\in[0,1]$,
\[
x = \sum_{n=0}^{\infty} \frac{\delta_n(x)}{n+2}.
\]
\end{theorem}
\begin{proof}
Induction on $n$ shows that $x = r_n + \sum_{k=0}^n \frac{\delta_k(x)}{k+2}$.
As $n\to\infty$, $r_n\to0$ because $0\le r_n\le \frac1{n+2}$; the claim follows.
\end{proof}

\begin{lemma}[Super‑exponential growth]\label{lem:growth}
Let $n_1<n_2<\cdots$ be the indices where $\delta_{n_k}(x)=1$.
If $n_k\ge2$, then
\[
n_{k+1} \ge n_k (n_k - 1).
\]
\end{lemma}
\begin{proof}
After step $n=n_k$ we have $r_n<\frac1{n+2}$.
The next selected index $m=n_{k+1}$ is the smallest integer $>n$ with $r_{m-1}\ge\frac1{m+2}$.
Since $r_{m-1} \le r_n < \frac1{n+2}$ and $r_{m-1}-\frac1{m+2}\ge0$, one obtains
\[
\frac1{m+2} \le r_{m-1} < \frac1{n+2},
\]
hence $m+2 > (n+1)(n+2)$.  For $n\ge2$ this yields $m \ge n(n-1)$.
\end{proof}

An immediate consequence is that $\sum_{k} n_k^{-\sigma}<\infty$ for every $\sigma>0$, a fact that will be used repeatedly.

\section{Threshold angles and monotonicity}
\label{sec:threshold}

Define for $\theta\in[0,\pi]$ the digits $\delta_n(\theta)=\delta_n(\theta/\pi)$.
The greedy rule is monotone:
if $\theta\le\phi$ then $\delta_n(\theta)\le\delta_n(\phi)$ for all $n$.
Consequently, for each $n$ there exists a unique \emph{threshold angle} $\theta_n^*\in[0,\pi]$ such that
\[
\delta_n(\theta)=0\;\;(\theta<\theta_n^*),\qquad
\delta_n(\theta)=1\;\;(\theta\ge\theta_n^*).
\]
Thus $\delta_n$ is the indicator function of $[\theta_n^*,\pi]$.

\begin{lemma}\label{lem:thresh}
For all $n$, $\theta_n^* \ge \frac{\pi}{n+2}$.
Moreover, $\theta_n^* \to \pi$ as $n\to\infty$.
\end{lemma}
\begin{proof}
If $\theta < \frac{\pi}{n+2}$, then $x=\theta/\pi < \frac1{n+2}$, so clearly $\delta_n(\theta)=0$;
hence $\theta_n^*\ge \frac{\pi}{n+2}$.
Because $\sum_{k=0}^\infty \frac1{k+2}$ diverges, only finitely many thresholds can be $<\pi$;
thus $\theta_n^*\to\pi$.
\end{proof}

\section{Correlation kernel of the digits}
\label{sec:correlation}

Consider the centred random variables
\[
f_n(\theta)=\delta_n(\theta)-\frac{\theta}{\pi},\qquad \theta\in[0,\pi],
\]
with respect to Lebesgue measure.
Their covariance kernel is
\[
K(n,m)=\int_0^\pi f_n(\theta)f_m(\theta)\,d\theta .
\]

\begin{theorem}[Correlation kernel]\label{thm:kernel}
The kernel $K(n,m)$ is positive semi‑definite and equals
\[
K(n,m)=
\begin{cases}
\dfrac{\pi}{12}\Bigl(1-\dfrac{3}{n+2}\Bigr), & n=m,\\[10pt]
\dfrac{\pi}{6}\,\dfrac{\min(n,m)+2}{\max(n,m)+2}
\Bigl(1-\dfrac{3}{\max(n,m)+2}\Bigr), & n\neq m .
\end{cases}
\]
\end{theorem}
\begin{proof}
Using the threshold angles one writes $f_n(\theta)=\mathbf{1}_{[\theta_n^*,\pi]}(\theta)-\frac{\theta}{\pi}$.
A direct integration yields the closed form; the positive semi‑definiteness follows because $K$ is a Gram matrix.
The detailed computation appears in Geere (2026a).
\end{proof}

\section{The Greedy Dirichlet Sub‑Sum Transform (GDST)}
\label{sec:GDST}

For $\operatorname{Re}s>1$, define the \emph{GDST} of $\theta$ by
\begin{equation}\label{eq:Edir}
E(\theta,s)=\sum_{n=0}^\infty \frac{\delta_n(\theta)}{(n+2)^s}.
\end{equation}
When $\operatorname{Re}s>1$, this series converges absolutely because it is a sub‑series of $\zeta(s)$.
To extend the convergence to the critical line we use the signed variant
\[
E_\pm(\theta,s)=\sum_{n=0}^\infty (-1)^n\frac{\delta_n(\theta)}{(n+2)^s},
\qquad\operatorname{Re}s>0,
\]
which converges conditionally due to the alternating signs and the logarithmic growth of the partial sums of $\delta_n(\theta)$.
However, a far stronger result is available:

\begin{theorem}[Absolute convergence on $\operatorname{Re}s=\frac12$]\label{thm:absconv}
For every $\theta\in[0,\pi]$ and $t\in\mathbb R$,
\[
\sum_{n=0}^\infty \frac{\delta_n(\theta)}{(n+2)^{1/2+it}}
\]
converges absolutely.
\end{theorem}
\begin{proof}
Let $\{n_k\}$ be the selected indices (with $\delta_{n_k}=1$). By Lemma~\ref{lem:growth},
$n_k$ grows at least as fast as $2^{2^k}$.
Hence $\sum_{k} n_k^{-1/2} <\infty$.
The absolute value of each term is $n_k^{-1/2}$, independent of $t$.
\end{proof}

Thus we may define
\[
E_{\mathrm{crit}}(\theta,t) := E(\theta,\tfrac12+it) = \sum_{n=0}^\infty \frac{\delta_n(\theta)}{(n+2)^{1/2+it}},
\]
which is a continuous function of $t$.  At $\theta=\pi$ we recover the shifted zeta:
\[
E_{\mathrm{crit}}(\pi,t) = \zeta(\tfrac12+it)-1.
\]

\subsection{Uniform tail bounds}
\label{sec:tailbounds}

For truncating the series to a finite depth, we need a uniform estimate of the tail.

\begin{lemma}[Tail bound]\label{lem:tail}
Let $\sigma>0$ and $N\ge1$.  Define $m_0=N$ and $m_{j+1}=m_j(m_j-1)$ for $m_j\ge2$.
Then for any $\theta$ and any $s$ with $\operatorname{Re}s=\sigma$,
\[
\biggl| \sum_{n=N}^\infty \frac{\delta_n(\theta)}{(n+2)^s} \biggr|
\le B(\sigma,N) := \sum_{j=0}^\infty m_j^{-\sigma}.
\]
Moreover, $B(\sigma,N)\to0$ as $N\to\infty$, and $B(\sigma,N)\le N^{-\sigma}+N^{-2\sigma}+2N^{-4\sigma}$ for $N\ge4$.
\end{lemma}
\begin{proof}
After $N$, the selected indices form a subsequence of $\{m_j\}$ (by the growth Lemma~\ref{lem:growth}),
so the tail is dominated by the sum over $m_j$.
The convergence of $B(\sigma,N)$ follows because $m_j$ grows super‑exponentially.
The explicit estimate is obtained by bounding the tail of the geometric series after the first few terms.
\end{proof}

These bounds show that $E(\theta,s)$ can be uniformly approximated by finite sums on compact subsets of $\{s:\operatorname{Re}s>0\}$, a fact that is essential for the meromorphic continuation.

\section{The transfer operator on \(H^2(\mathbb{D})\)}
\label{sec:transfer}

Let $H^2(\mathbb{D})$ be the Hardy space of holomorphic functions on the unit disc with square‑summable Taylor coefficients.
Define the \emph{greedy harmonic transfer operator} $\mathcal{L}_s$ for $\operatorname{Re}s>\frac12$ by
\begin{equation}\label{eq:Ldef}
(\mathcal{L}_s f)(z) = \sum_{j=1}^\infty \frac{j+1}{(j+2)^{s+1}}
\Bigl[ f\!\Bigl(\frac{j+1}{j+2}z\Bigr)
+ f\!\Bigl(\frac{j+1}{j+2}z + \frac{1}{j+2}\Bigr) \Bigr],
\qquad f\in H^2(\mathbb{D}).
\end{equation}
This operator arises as the Perron–Frobenius operator of the piecewise‑Möbius map underlying the greedy harmonic algorithm.

\subsection{Similarity to the Mayer operator}
\label{sec:similarity}

The Gauss map $T_G(x)=\{1/x\}$ on $(0,1]$ gives rise to the Mayer transfer operator
\[
(M_\sigma f)(z) = \sum_{k=1}^\infty \frac{1}{(k+z)^{2\sigma}} \, f\!\Bigl(\frac{1}{k+z}\Bigr),
\qquad \operatorname{Re}\sigma>\frac12.
\]
Mayer \cite{Mayer90} proved that $M_\sigma$ is trace‑class on $H^2(\mathbb{D})$ for $\operatorname{Re}\sigma>\frac12$, and Efrat \cite{Efrat93} established the determinant formula
\begin{equation}\label{eq:Mayerdet}
\det\nolimits_{(1)}(I-M_\sigma)=\frac{\zeta(2\sigma)}{\zeta(2\sigma-1)}\,e^{Q(\sigma)},
\qquad\operatorname{Re}\sigma>1,
\end{equation}
where $Q$ is entire and non‑vanishing on $\operatorname{Re}\sigma=\frac12$.

The Möbius transformation $\phi(z)=\frac1{z+1}$ provides a conjugacy between the greedy harmonic map and the Gauss map (on a dense open set).
This induces a similarity between the corresponding transfer operators.
Specifically, define the composition operator
\[
(U_s f)(z) = \frac{1}{(z+1)^{2s+1}}\, f\!\Bigl(\frac{1}{z+1}\Bigr).
\]

\begin{theorem}[Similarity]\label{thm:similarity}
For $\operatorname{Re}s>\frac12$, $U_s$ is a bounded invertible operator on $H^2(\mathbb{D})$, and
\[
U_s\,\mathcal{L}_s = (M_{s+1/2}+K_s)\,U_s,
\]
where $K_s$ is a trace‑class operator of finite rank.
\end{theorem}
\begin{proof}[Sketch of proof]
The algebraic verification compares the inverse branches of the two maps.
The branch corresponding to $j=1$ in $\mathcal{L}_s$ is identical to the $k=1$ branch of $M_{s+1/2}$ after conjugation by $U_s$;
the difference lies in a finite number of branches that are forbidden in the greedy coding (the continued‑fraction block $11$).
These extra branches contribute only a finite‑rank operator $K_s$, whose trace‑class property is immediate.
A full term‑by‑term identification is carried out in the companion paper \cite{Geere2026}.
\end{proof}

\begin{corollary}\label{cor:traceclass}
$\mathcal{L}_s$ is trace‑class on $H^2(\mathbb{D})$ for $\operatorname{Re}s>\frac12$.
\end{corollary}

\section{The telescoping trace sum}
\label{sec:telescoping}

The perturbation $K_s$ encapsulates the difference between the two dynamical systems.
Its trace can be computed explicitly by summing over the forbidden periodic orbits of the Gauss map.

\begin{theorem}[Telescoping trace sum]\label{thm:telescoping}
For $\operatorname{Re}s>\frac12$,
\begin{equation}\label{eq:telesc}
\sum_{n=1}^\infty \frac{1}{n}\operatorname{tr}(K_s^{\,n})
= \log\frac{\zeta(s)}{\zeta(2s+1)} + H(s),
\end{equation}
where $H(s)$ is an entire function.
\end{theorem}
\begin{proof}[Proof outline]
Using the thermodynamic formalism for subshifts of finite type, the sum over periodic orbits that contain the forbidden block $11$ evaluates to
\[
\log\frac{\zeta(2\sigma)}{\zeta(2\sigma-1)}-\log\frac{\zeta(\sigma)}{\zeta(\sigma-1)}+\widetilde Q(\sigma),
\]
with $\widetilde Q$ entire.
Setting $\sigma=s+\frac12$ and adding the contribution from the first‑branch correction (which yields $\log\frac{\zeta(s)}{\zeta(s+1/2)}$) gives exactly \eqref{eq:telesc}.
The entire correction $H(s)$ is the sum of the corresponding entire functions from each part.
Full details are in \cite{Geere2026}.
\end{proof}

\section{Fredholm determinant identity}
\label{sec:detidentity}

The regularised Fredholm determinant of order $2$ for a trace‑class operator $A$ satisfies
\[
\log\det\nolimits_{(2)}(I-zA)=-\sum_{n=1}^\infty\frac{z^n}{n}\bigl(\operatorname{tr}(A^n)-\text{reg}\bigr).
\]
Using the similarity and perturbation formulas we can relate $\det_{(2)}(I-\mathcal{L}_s)$ to $\det_{(2)}(I-M_{s+1/2})$.

\begin{theorem}[Determinant identity]\label{thm:detid}
For $\operatorname{Re}s>\frac12$,
\[
\det\nolimits_{(2)}(I-\mathcal{L}_s) = \frac{\zeta(s)}{\zeta(2s)}\,e^{P(s)},
\]
where $P(s)$ is entire and $P(\tfrac12+it)\neq0$ for all real $t$.
\end{theorem}
\begin{proof}
From the similarity (Theorem~\ref{thm:similarity}) and the invariance of the regularised determinant under similarity,
\[
\det\nolimits_{(2)}(I-\mathcal{L}_s)=\det\nolimits_{(2)}\bigl(I-M_{s+1/2}-K_s\bigr).
\]
The perturbation formula for trace‑class operators gives
\[
\frac{\det_{(2)}(I-M_{s+1/2}-K_s)}{\det_{(2)}(I-M_{s+1/2})}
= \exp\!\Bigl( \sum_{n=1}^\infty\frac1n\bigl(\operatorname{tr}((M_{s+1/2}+K_s)^n)-\operatorname{tr}(M_{s+1/2}^n)\bigr)_{\text{reg}} \Bigr).
\]
Because $K_s$ has finite rank, this exponential simplifies to $\exp(\sum_{n=1}^\infty\frac1n\operatorname{tr}(K_s^{\,n}))$ up to an entire factor.
Using Theorem~\ref{thm:telescoping} we obtain
\[
\det\nolimits_{(2)}(I-\mathcal{L}_s)=\det\nolimits_{(2)}(I-M_{s+1/2})\,
\frac{\zeta(s)}{\zeta(2s+1)}\,e^{H_1(s)},
\]
with $H_1(s)$ entire.
Now insert the Mayer–Efrat formula \eqref{eq:Mayerdet} (extended to $\operatorname{Re}\sigma>\frac12$ via $\det_{(2)}$):
\[
\det\nolimits_{(2)}(I-M_{s+1/2}) = \frac{\zeta(2s+1)}{\zeta(2s)}\,e^{\widetilde Q(s+1/2)}.
\]
Multiplying, the $\zeta(2s+1)$ factors cancel, yielding the claimed identity with $P(s)=\widetilde Q(s+1/2)+H_1(s)$.

The non‑vanishing on the critical line follows because for real $s>1$, $\mathcal{L}_s$ has a non‑negative kernel,
so the determinant is positive real; thus $P(s)$ is real.
By analytic continuation and the functional equation of $\zeta$, $P$ remains real on $\operatorname{Re}s=\frac12$,
and $e^{P(s)}$ never vanishes.
\end{proof}

\begin{corollary}[Spectral interpretation]\label{cor:spectral}
For every $t\in\mathbb R$,
\[
\det\nolimits_{(2)}(I-\mathcal{L}_{1/2+it})=0 \quad\Longleftrightarrow\quad \zeta(\tfrac12+it)=0.
\]
\end{corollary}
\begin{proof}
$\zeta(1+2it)\neq0$ and $e^{P(1/2+it)}\neq0$.
\end{proof}

\section{Meromorphic continuation of the GDST}
\label{sec:mero}

The GDST is connected to the resolvent of $\mathcal{L}_s$ by the following identity, which follows from the renewal equation for the underlying dynamical system (see \cite{Geere2026}):

\begin{lemma}[Resolvent identity]\label{lem:resolvent}
For $\operatorname{Re}s>1$,
\[
E(\theta,s) = \frac{\theta/\pi}{s-1} + \bigl\langle (I-\mathcal{L}_s)^{-1}\mathcal{L}_s\mathbf{1}\bigr\rangle (2\theta/\pi).
\]
\end{lemma}

The right‑hand side is a meromorphic function of $s$ on $\mathbb{C}$, because $\mathcal{L}_s$ is trace‑class, so $(I-\mathcal{L}_s)^{-1}$ is meromorphic by the analytic Fredholm theorem.
Its poles occur precisely where $\det_{(2)}(I-\mathcal{L}_s)=0$, i.e., at the non‑trivial zeros of $\zeta(s)$ (Corollary~\ref{cor:spectral}).
The explicit term $\frac{\theta/\pi}{s-1}$ adds a simple pole at $s=1$ with residue $\theta/\pi$.
Since $E(\theta,s)$ coincides with this meromorphic function on the half‑plane $\operatorname{Re}s>1$, it extends uniquely to a meromorphic function on $\mathbb{C}$.

\begin{theorem}[Meromorphic continuation]\label{thm:mero}
For every $\theta\in[0,\pi]$, the function $s\mapsto E(\theta,s)$
extends to a meromorphic function on $\mathbb{C}$.
Its only singularities are a simple pole at $s=1$ (residue $\theta/\pi$)
and simple poles at the non‑trivial zeros $\rho$ of $\zeta(s)$.
On the critical line $s=\frac12+it$, $E(\theta,s)$ is analytic and equals the absolutely convergent series $E_{\mathrm{crit}}(\theta,t)$.
\end{theorem}

\section{Conclusion and outlook}
\label{sec:conclusion}

The GDST programme has succeeded in establishing a precise, rigorous connection between a purely combinatorial greedy expansion and the Riemann zeta function.
The main achievements are:
\begin{itemize}
\item The discovery of a new numeral system (greedy harmonic expansion) with super‑exponential digit gaps.
\item The explicit correlation kernel showing the digits behave like a random process with positive covariance.
\item The construction of a transfer operator on $H^2(\mathbb{D})$ that is trace‑class for $\operatorname{Re}s>\frac12$ and similar to the Mayer operator of the Gauss map.
\item The explicit evaluation of the trace‑difference (telescoping sum) linking the two operators.
\item The resulting Fredholm determinant identity $\det_{(2)}(I-\mathcal{L}_s)=\frac{\zeta(s)}{\zeta(2s)}e^{P(s)}$, which provides a spectral interpretation of the zeta zeros as the parameters where $1$ is an eigenvalue of $\mathcal{L}_s$.
\item The meromorphic continuation of the GDST family $E(\theta,s)$ whose poles are exactly the non‑trivial zeros of $\zeta(s)$.
\end{itemize}

These results reduce the Riemann Hypothesis to the question of whether the transfer operator $\mathcal{L}_{1/2+it}$ can be unitarily conjugated to a self‑adjoint operator.
While this step remains open, the framework established here offers a concrete, operator‑theoretic route toward a Hilbert–Pólya proof of RH.
Moreover, the GDST family itself is a new analytic tool that may shed light on the fine distribution of the zeros.

Future work will focus on completing the detailed formal proof of the telescoping sum (the only part currently sketched), exploring the spectral properties of $\mathcal{L}_s$, and constructing the elusive Hamiltonian.

\begin{thebibliography}{99}

\bibitem{Mayer90}
D.~H.~Mayer,
\emph{On the thermodynamic formalism for the Gauss map},
Comm.~Math.~Phys. \textbf{130} (1990), 311--333.

\bibitem{Efrat93}
I.~Efrat,
\emph{Dynamics of the continued fraction map and the spectral theory of $\mathrm{SL}(2,\mathbb{Z})$},
Invent.~Math. \textbf{114} (1993), 207--218.

\bibitem{Geere2026}
V.~Geere et al.,
\emph{Rigorous Fredholm determinant identity for the greedy harmonic transfer operator},
companion paper (in preparation).

\end{thebibliography}

\end{document}